\documentclass[11pt,a4paper,oneside]{report}
\usepackage[spanish,activeacute,es-noshorthands]{babel}
\usepackage[utf8]{inputenc}
\usepackage[T1]{fontenc}
\usepackage{lmodern}
\usepackage{amsmath, amssymb}
\usepackage{graphicx}
\usepackage[dvipsnames]{xcolor}
\usepackage{geometry}
\geometry{margin=2.4cm, headheight=14pt}
\usepackage{microtype}
\usepackage{booktabs}
\usepackage{longtable}
\usepackage{enumitem}
\usepackage{fancyhdr}
\usepackage{tcolorbox}
\tcbuselibrary{skins, breakable}
\usepackage{caption}
\captionsetup{font=small, labelfont=bf}
\usepackage{titlesec}
\usepackage{hyperref}
\hypersetup{
  colorlinks=true,
  linkcolor=NavyBlue,
  citecolor=ForestGreen,
  urlcolor=BrickRed,
  pdftitle={Guia de empleo en automatizacion industrial y mineria - Peru},
  pdfauthor={CubeSat-EdgeAI-Payload}
}

\definecolor{intisat}{RGB}{26, 54, 93}
\definecolor{intisataccent}{RGB}{197, 143, 0}
\definecolor{ok}{RGB}{47, 133, 90}
\definecolor{warn}{RGB}{197, 60, 60}

\titleformat{\chapter}[hang]
  {\normalfont\Huge\bfseries\color{intisat}}
  {\thechapter.}{0.6em}{}
\titlespacing*{\chapter}{0pt}{-20pt}{20pt}

\pagestyle{fancy}
\fancyhf{}
\fancyhead[L]{\nouppercase{\leftmark}}
\fancyhead[R]{\thepage}
\renewcommand{\headrulewidth}{0.4pt}

\newtcolorbox{tip}[1][]{
  colback=green!5, colframe=ok,
  fonttitle=\bfseries, title=Tip practico,
  breakable, sharp corners=south, #1
}
\newtcolorbox{advertencia}[1][]{
  colback=red!4, colframe=warn,
  fonttitle=\bfseries, title=Advertencia,
  breakable, sharp corners=south, #1
}
\newtcolorbox{datos}[1][]{
  colback=blue!4, colframe=intisat,
  fonttitle=\bfseries, title=Datos del mercado,
  breakable, sharp corners=south, #1
}

\begin{document}

% ─── Portada ────────────────────────────────────────────────────────
\begin{titlepage}
\centering
\vspace*{1cm}
{\color{intisat}\rule{\textwidth}{2pt}}
\vspace{0.6cm}

{\Huge\bfseries\color{intisat} Guia de empleo}\\[0.3em]
{\LARGE\bfseries Automatizacion industrial y mineria}\\[0.3em]
{\Large Peru — 2026}

\vspace{0.7cm}
{\Large\itshape Sistemas industriales, sueldos, empresas\\
y plan de accion para transicionar desde\\
un proyecto aeroespacial academico al sector productivo}

\vspace{1cm}
{\color{intisataccent}\rule{0.6\textwidth}{1pt}}

\vspace{2cm}

\begin{tabular}{rl}
\textbf{Audiencia:} & Estudiantes / egresados de ingenieria \\
\textbf{Sectores:} & Mineria, alimentos, manufactura \\
\textbf{Roles:} & Automatizacion, control, data science industrial \\
\textbf{Origen:} & Proyecto CubeSat-EdgeAI-Payload (UNI / INTISAT) \\
\textbf{Fecha:} & \today \\
\textbf{Documento:} & v1.0
\end{tabular}

\vfill

\begin{minipage}{0.85\textwidth}
\textbf{Resumen.}\quad
Documento orientativo sobre el sector de automatizacion industrial en
Peru, con foco en mineria, alimentos y manufactura. Incluye descripcion
de sistemas (PLC, SCADA, DCS), buses de comunicacion, marcas dominantes,
sueldos reales por nivel de experiencia, empresas tipicas, plataformas
de busqueda de empleo, certificaciones recomendadas, y un plan de
transicion concreto para egresados de proyectos aeroespaciales /
academicos hacia el sector industrial productivo. Las cifras son
referenciales y deben verificarse contra las plataformas indicadas.
\end{minipage}

\vspace{1.2cm}
{\color{intisat}\rule{\textwidth}{2pt}}
\end{titlepage}

\tableofcontents
\clearpage

% =====================================================================
\chapter{Vision general del sector}

Peru tiene una base industrial significativa en tres sectores:

\begin{itemize}[leftmargin=2em]
  \item \textbf{Mineria}: la columna vertebral. Operaciones de gran escala
        (Antamina, Yanacocha, Cerro Verde, Las Bambas, Toromocho, Quellaveco)
        usan automatizacion identica a Australia, Canada o Chile. Inversion
        anual en sistemas de control: cientos de millones de USD.
  \item \textbf{Alimentos y bebidas}: empresas grandes (Alicorp, Gloria,
        Backus/AB InBev, Nestle Peru, Laive) tienen lineas automatizadas
        con PLCs, MES, trazabilidad batch-to-end.
  \item \textbf{Manufactura}: mas chico que mineria/alimentos pero existe
        (Aceros Arequipa, Continental Tires, Lindley, Ferreyros).
\end{itemize}

\begin{datos}
La inversion total en sistemas de control industrial en Peru se estima
en USD 500--800 millones por año. La mayor parte es renovacion / expansion
de plantas existentes en mineria y alimentos.
\end{datos}

\section{Donde NO estamos atrasados}

A diferencia del sector aeroespacial (donde Peru tiene poca capacidad
local), en industrial:

\begin{itemize}[leftmargin=2em]
  \item Las grandes mineras usan exactamente el mismo hardware
        (Siemens, Allen-Bradley, ABB) que sus operaciones globales.
  \item Hay representantes oficiales y stock local de marcas premium.
  \item Existen tecnicos peruanos con certificaciones internacionales.
  \item Los integradores locales pueden hacer proyectos llave en mano.
\end{itemize}

\section{Donde si hay limitaciones}

\begin{itemize}[leftmargin=2em]
  \item Pocas empresas peruanas \emph{fabrican} hardware de control;
        casi todo se importa.
  \item La capacidad de R\&D propia es limitada.
  \item Brain drain hacia Chile / Australia / Canada en niveles senior.
  \item Salaries en automatizacion son altos en mineria pero bajos en
        otros sectores.
\end{itemize}

% =====================================================================
\chapter{Sistemas industriales: PLC, DCS, SCADA, MES}

\section{PLC --- Programmable Logic Controller}

Computadora especializada en automatizacion. A diferencia de un PC, esta
diseñada para:

\begin{itemize}[leftmargin=2em]
  \item Operar 24/7/365 durante 20+ años sin reboot.
  \item Soportar -40 a +70 °C en sala industrial polvorienta.
  \item Resistir ruido electrico de motores y soldadoras.
  \item Garantizar tiempo de respuesta deterministico.
  \item Programarse con lenguajes pensados para procesos.
\end{itemize}

\subsection{Marcas dominantes mundialmente}

\begin{center}
\begin{tabular}{lll}
\toprule
\textbf{Marca} & \textbf{Producto estrella} & \textbf{Donde domina} \\
\midrule
Siemens               & SIMATIC S7-1200 / 1500   & mineria europea/peruana, refinerias \\
Rockwell Allen-Bradley & ControlLogix / CompactLogix & USA, comida, automotive \\
Schneider Electric    & Modicon M340 / M580      & Francia, energia \\
Mitsubishi            & MELSEC iQ-R              & Japon, automotive \\
Omron                 & Sysmac NX/NJ             & electronica, packaging \\
ABB                   & AC500                    & mineria, energia \\
Beckhoff              & TwinCAT (PC-based)       & maquinaria europea high-end \\
Honeywell             & Experion PKS (DCS)       & refinerias, oil \& gas \\
\bottomrule
\end{tabular}
\end{center}

En Peru predominan \textbf{Siemens} (mineria) y \textbf{Allen-Bradley}
(alimentos como Backus, Alicorp).

\section{DCS --- Distributed Control System}

Variante del PLC para procesos continuos grandes (refinerias, plantas
quimicas, plantas concentradoras de mineral). La diferencia clave:

\begin{itemize}[leftmargin=2em]
  \item PLC: discreto, eventos rapidos, una maquina o linea.
  \item DCS: continuo, miles de loops PID, planta entera coordinada.
\end{itemize}

Marcas DCS: Honeywell Experion PKS, Yokogawa CENTUM VP, ABB 800xA,
Emerson DeltaV.

\section{SCADA --- Supervisory Control And Data Acquisition}

El \emph{software} que el operador ve en la pantalla. Conecta a los PLCs/DCS
y muestra el proceso en tiempo real con sinopticos, alarmas, tendencias,
historicos.

\begin{center}
\begin{tabular}{ll}
\toprule
\textbf{Plataforma SCADA} & \textbf{Quien la fabrica} \\
\midrule
AVEVA InTouch (ex-Wonderware) & AVEVA / Schneider \\
Ignition                       & Inductive Automation \\
WinCC                         & Siemens \\
iFix                          & GE Digital \\
Citect SCADA                  & Schneider Electric \\
ClearSCADA                    & AVEVA \\
\bottomrule
\end{tabular}
\end{center}

En mineria peruana es muy comun ver \textbf{Wonderware InTouch} y
\textbf{Ignition}. \textbf{WinCC} en plantas Siemens-centric.

\section{MES --- Manufacturing Execution System}

Capa por encima del SCADA. Conecta el piso de planta con la oficina
(ERP). Maneja:

\begin{itemize}[leftmargin=2em]
  \item Trazabilidad batch-to-end (que materia prima fue a que producto).
  \item OEE (Overall Equipment Effectiveness) y KPI de planta.
  \item Calidad y validacion (FDA / DIGESA en alimentos).
  \item Programacion de produccion en tiempo real.
\end{itemize}

Marcas: AVEVA MES, Siemens Opcenter, Rockwell PharmaSuite, GE Proficy.

\section{Arquitectura ISA-95 (la piramide de automatizacion)}

\begin{center}
\begin{tabular}{lll}
\toprule
\textbf{Nivel} & \textbf{Que es} & \textbf{Tipico} \\
\midrule
Nivel 5 & ERP                   & SAP, Oracle. Compras, RR.HH., finanzas \\
Nivel 4 & MES / Planificacion    & AVEVA MES, OEE, calidad \\
Nivel 3 & SCADA / Operacion      & Ignition, Wonderware, WinCC \\
Nivel 2 & Control                & PLC, DCS \\
Nivel 1 & Instrumentacion        & Sensores, actuadores, drivers \\
Nivel 0 & Proceso                & Materia prima, horno, mina \\
\bottomrule
\end{tabular}
\end{center}

% =====================================================================
\chapter{Buses y protocolos industriales}

A diferencia de I2C/SPI/UART (que son intra-PCB, distancias cortas), los
buses industriales son \emph{deterministic, isolation transformer, longitud
larga, ruido electrico tolerante}.

\begin{center}
\begin{tabular}{llll}
\toprule
\textbf{Protocolo} & \textbf{Distancia} & \textbf{Velocidad} & \textbf{Uso} \\
\midrule
4-20 mA (analogico) & 1 km     & ---       & sensores presion, temp, nivel \\
HART                & 1 km     & 1.2 kbps  & 4-20 mA + datos digitales \\
Modbus RTU (RS-485) & 1 km     & 115 kbps  & el "I2C industrial" \\
Modbus TCP          & 100 m    & 100 Mbps  & version IP de Modbus \\
Profibus DP         & 100 m    & 12 Mbps   & red Siemens clasica \\
Profinet            & 100 m    & 1 Gbps    & Siemens IP-based moderno \\
EtherCAT            & 100 m    & 100 Mbps  & motion control microsegundo \\
DeviceNet / CIP     & 500 m    & 500 kbps  & Allen-Bradley legacy \\
OPC UA              & red IP   & ---       & moderno, vendor-neutral \\
ISA-100 / WirelessHART & inalambrico & --- & sensores remotos \\
\bottomrule
\end{tabular}
\end{center}

\textbf{El gran cambio reciente} es \textbf{OPC UA}: protocolo abierto,
vendor-neutral, que permite interoperabilidad sin lock-in propietario.

\section{Lenguajes de programacion (norma IEC 61131-3)}

\begin{center}
\begin{tabular}{ll}
\toprule
\textbf{Lenguaje} & \textbf{Uso} \\
\midrule
Ladder Diagram (LD)        & logica binaria, releve emulado. El mas usado. \\
Function Block Diagram (FBD) & flujo de datos, control PID \\
Structured Text (ST)        & algoritmos complejos (parecido a Pascal) \\
Sequential Function Chart (SFC) & maquinas de estado \\
Instruction List (IL)       & assembler-like, casi extinto \\
\bottomrule
\end{tabular}
\end{center}

% =====================================================================
\chapter{Sectores en detalle: mineria}

La mineria peruana es donde \textbf{mas se invierte} en automatizacion. Las
grandes operaciones tienen sistemas de primer nivel mundial.

\section{Empresas operadoras}

\begin{center}
\begin{tabular}{p{4.5cm}p{4cm}p{4cm}}
\toprule
\textbf{Empresa} & \textbf{Operacion} & \textbf{Ubicacion} \\
\midrule
Antamina (BHP/Glencore/Teck/Mitsubishi)  & cobre/zinc & Ancash (4 200 m) \\
Yanacocha (Newmont)                      & oro        & Cajamarca \\
Cerro Verde (Freeport-McMoRan)           & cobre      & Arequipa \\
Las Bambas (MMG)                         & cobre      & Apurimac (4 000 m) \\
Toromocho (Chinalco)                     & cobre      & Junin (4 500 m) \\
Quellaveco (Anglo American)              & cobre      & Moquegua \\
Cuajone (Southern Copper)                & cobre      & Moquegua \\
Volcan                                   & polimetalico & Cerro de Pasco \\
Buenaventura                             & polimetalico & varias \\
Hochschild Mining                        & oro/plata    & varias \\
\bottomrule
\end{tabular}
\end{center}

\section{Sistemas tipicos por proceso}

\begin{center}
\begin{tabular}{p{3.5cm}p{4.5cm}p{4.5cm}}
\toprule
\textbf{Proceso} & \textbf{Tecnologia} & \textbf{Marca tipica} \\
\midrule
Trituracion primaria & PLC + variadores & Siemens S7-400 + Sinamics \\
Molinos SAG/bolas    & DCS o PLC redundante & Honeywell Experion / ABB 800xA \\
Flotacion            & analizador en linea + PLC & Outotec Courier + Siemens \\
Espesadores          & sensores nucleares + PLC & Endress+Hauser + Allen-Bradley \\
Camiones autonomos   & sistema completo & Komatsu FrontRunner / Cat MineStar \\
Despacho             & software optimizacion & Modular Mining DISPATCH \\
SCADA central        & sinopticos + alarmas & Wonderware / Ignition \\
\bottomrule
\end{tabular}
\end{center}

\begin{datos}
\textbf{Inversion tipica}: una planta concentradora de 100 000 toneladas/dia
gasta entre USD 30 y 80 millones \textbf{solo en sistema de control}, sin
contar instrumentacion. El presupuesto de automatizacion es 5--10\% del
total de la planta.
\end{datos}

% =====================================================================
\chapter{Sectores en detalle: alimentos y bebidas}

\section{Empresas grandes}

\begin{center}
\begin{tabular}{p{3.5cm}p{4cm}p{4cm}}
\toprule
\textbf{Empresa} & \textbf{Sector} & \textbf{Automatizacion tipica} \\
\midrule
Alicorp                       & aceites, fideos, detergentes & Allen-Bradley + RSLogix \\
Gloria                        & lacteos & Siemens + WinCC \\
Backus (AB InBev)             & cerveza & Siemens TIA Portal + WinCC \\
Laive                         & lacteos & mix \\
Nestle Peru                   & lacteos, cafe & global standard ABB / Siemens \\
San Fernando                  & avicola & Siemens + WinCC \\
Lindley (Coca-Cola Peru)      & bebidas & Siemens \\
Aje                          & bebidas & Allen-Bradley \\
\bottomrule
\end{tabular}
\end{center}

\section{Caracteristicas tecnicas}

\begin{itemize}[leftmargin=2em]
  \item \textbf{Esterilizacion / pasteurizacion} controlada por PLC con
        curvas de temperatura validadas (FDA / DIGESA).
  \item \textbf{Llenado y empaque} con servomotores Siemens / Beckhoff
        sincronizados a milisegundo.
  \item \textbf{Trazabilidad batch-to-end} con MES (Wonderware MES, Siemens
        Opcenter) que conecta con ERP SAP.
  \item \textbf{Vision por computadora} para inspeccion de envases (Cognex,
        Keyence). Aqui es donde tu skill IA tiene aplicacion directa.
\end{itemize}

% =====================================================================
\chapter{Sectores en detalle: manufactura}

Mas chico que mineria/alimentos pero existe:

\begin{center}
\begin{tabular}{lll}
\toprule
\textbf{Empresa} & \textbf{Sector} & \textbf{Sistema} \\
\midrule
Aceros Arequipa       & siderurgia        & Siemens + ABB drives \\
Ferreyros (CAT)       & maquinaria pesada & Allen-Bradley \\
Continental Tires     & neumaticos        & Mitsubishi \\
Cogorno               & textiles          & mix economico \\
Topitop               & textiles          & mix \\
Sider Peru            & acero             & Siemens \\
Yura (Cementos Yura)  & cemento           & ABB \\
UNACEM                & cemento           & ABB / Siemens \\
\bottomrule
\end{tabular}
\end{center}

% =====================================================================
\chapter{Sueldos por nivel de experiencia}

\begin{advertencia}
Cifras referenciales basadas en publicaciones del sector y reportes de
LinkedIn / Bumeran / Computrabajo. Verificar siempre contra ofertas
actuales en las plataformas.
\end{advertencia}

Peru paga \textbf{14 sueldos al año} (12 mensuales + 2 gratificaciones de
julio y diciembre). Las cifras siguientes son sueldos brutos mensuales.

\begin{center}
\begin{tabular}{lp{2cm}p{2.6cm}p{2.6cm}p{3cm}}
\toprule
\textbf{Nivel} & \textbf{Exp. (años)} & \textbf{S/ mensual} & \textbf{USD mensual} & \textbf{USD anual (14)} \\
\midrule
Practicante / Trainee   & 0    & 1 500 - 2 500   & 400 - 680     & 5 600 - 9 500 \\
Junior                  & 1-3  & 4 000 - 7 000   & 1 100 - 1 900 & 15 000 - 26 000 \\
Semi-senior             & 3-7  & 7 000 - 12 000  & 1 900 - 3 300 & 26 000 - 46 000 \\
Senior                  & 7-15 & 12 000 - 22 000 & 3 300 - 6 000 & 46 000 - 84 000 \\
Lead / Especialista     & 15+  & 20 000 - 40 000 & 5 500 - 10 800 & 77 000 - 150 000 \\
Gerente Automatizacion  & 20+  & 35 000 - 70 000 & 9 500 - 18 900 & 132 000 - 265 000 \\
\bottomrule
\end{tabular}
\end{center}

Estas cifras \textbf{NO incluyen bonos}. La industria minera paga ademas:

\begin{itemize}[leftmargin=2em]
  \item \textbf{Bono anual por desempeño}: 1--3 sueldos extra.
  \item \textbf{Utilidades} (ley peruana): hasta 18 sueldos en años buenos.
  \item \textbf{Bono de mina} (si trabajas en site): 30--50\% adicional al base.
\end{itemize}

\textbf{Senior en site con buena empresa}: facilmente USD 150 000/año todo
incluido.

\section{Factores que cambian el numero}

\subsection{Site (mina) vs Lima (oficina)}

\begin{center}
\begin{tabular}{ll}
\toprule
\textbf{Modalidad} & \textbf{Diferencia} \\
\midrule
Site (mina) turno 14$\times$7 o 21$\times$7 & base + 30--50\% \\
Lima oficina                                & base \\
Hibrido (Lima + visitas a site)             & base + 10--20\% \\
\bottomrule
\end{tabular}
\end{center}

\subsection{Empresa (tier)}

\begin{center}
\begin{tabular}{lll}
\toprule
\textbf{Tier} & \textbf{Empresas} & \textbf{Senior tipico (S/)} \\
\midrule
Tier 1 (multinacional) & Antamina, Cerro Verde, Yanacocha, Las Bambas, Quellaveco & 20 000 - 25 000 \\
Tier 2 (mid-size)      & Volcan, Buenaventura, Hochschild, Southern Copper       & 14 000 - 18 000 \\
Tier 3 (junior)        & Sociedad Minera, Aruntani                                & 10 000 - 13 000 \\
EPC / contratistas     & Bechtel, Fluor, Hatch, Ausenco                          & 15 000 - 22 000 \\
Integrador local       & Comesa, Tetra Tech, GTD                                  &  9 000 - 15 000 \\
\bottomrule
\end{tabular}
\end{center}

\subsection{Especialidades premium (+20--40\% sobre base)}

\begin{itemize}[leftmargin=2em]
  \item Ciberseguridad industrial OT (escasez de talento).
  \item MES / Sistemas (Wonderware, Ignition).
  \item DCS especialista (Honeywell Experion, Yokogawa CENTUM).
  \item PLC programmer Senior con Allen-Bradley.
  \item Despacho minero (Modular Mining DISPATCH, Hexagon, Caterpillar).
  \item \textbf{Vision por computadora industrial} (donde tu skill aplica).
  \item Robotica de proceso.
\end{itemize}

\subsection{Idioma}

\begin{itemize}[leftmargin=2em]
  \item Ingles fluido: $+20\%$ en sueldo (multinacionales lo exigen).
  \item Portugues / mandarin: bonus extra pero raro.
\end{itemize}

\section{Comparativa con otras especialidades}

\begin{center}
\begin{tabular}{ll}
\toprule
\textbf{Especialidad} & \textbf{Senior tipico (S/ mensual)} \\
\midrule
Automatizacion mineria               & 12 000 - 22 000 \\
Procesos / Metalurgia mineria        & 15 000 - 25 000 \\
Geomecanica mineria                  & 12 000 - 20 000 \\
Petroleo / refineria                 & 14 000 - 28 000 \\
Software Backend Senior              & 10 000 - 25 000 \\
Software FAANG remoto USA            & 25 000 - 80 000 \\
Aeroespacial Peru (CONIDA, UNI)      &  5 000 - 12 000 \\
Investigador universitario senior    &  8 000 - 15 000 \\
Ingenieria civil construccion        &  8 000 - 15 000 \\
Electronica consumer                 &  5 000 - 10 000 \\
\bottomrule
\end{tabular}
\end{center}

\begin{datos}
\textbf{Hecho clave}: en Peru, mineria paga mejor que aeroespacial por un
factor de 2--3x. El talento se va a mineria. Si vienes de un proyecto
aeroespacial academico, transicionar al sector industrial es \textbf{economicamente
ventajoso} y tecnicamente accesible.
\end{datos}

% =====================================================================
\chapter{Camino tipico de carrera}

\begin{center}
\begin{tabular}{ll}
\toprule
\textbf{Año} & \textbf{Hito} \\
\midrule
0   & Practicante en EPC o integrador, 2 500 S/ \\
2   & Junior en EPC, 6 500 S/ \\
4   & Junior en mina (site), 9 000 S/ + bonos \\
7   & Semi-senior en mina, 14 000 S/ + bonos + utilidades \\
12  & Senior especialista (DCS / cyber / MES), 22 000 S/ + bonos \\
18  & Lead / consultor independiente, 30 000+ S/ o \$1500/dia \\
\bottomrule
\end{tabular}
\end{center}

\section{Comparativa internacional}

\begin{center}
\begin{tabular}{ll}
\toprule
\textbf{Pais} & \textbf{Senior automatizacion mineria (USD/año)} \\
\midrule
Peru      & 46 000 - 84 000 \\
Chile     & 70 000 - 120 000 \\
Australia (FIFO) & 130 000 - 250 000 \\
Canada    & 110 000 - 180 000 \\
USA (oil \& gas) & 100 000 - 200 000 \\
\bottomrule
\end{tabular}
\end{center}

\textbf{Australia} es el destino estrella para mineros peruanos: dobla o
triplica el sueldo, modalidad FIFO (Fly-In Fly-Out, 14 dias en mina + 14 en
casa).

% =====================================================================
\chapter{Empresas relevantes para tu perfil}

\begin{advertencia}
Las empresas listadas tienen actividades de innovacion / data science /
vision computacional segun publicaciones publicas, pero los programas
especificos cambian seguido. \textbf{Verificar siempre} contra LinkedIn
y las paginas oficiales antes de aplicar.
\end{advertencia}

\section{Mineras operadoras con foco en automatizacion / AI}

\begin{center}
\begin{tabular}{p{4cm}p{8cm}}
\toprule
\textbf{Empresa} & \textbf{Por que mira tu perfil} \\
\midrule
Antamina        & area de innovacion / data, sistemas de despacho avanzados \\
Cerro Verde     & operacion grande con renovacion tecnologica continua \\
Anglo American Quellaveco & "operacion 100\% digital" segun publicaciones; alto perfil de data \\
Las Bambas      & sistemas de despacho autonomos en evaluacion \\
Yanacocha       & area de transformacion digital activa \\
Buenaventura    & varias operaciones, oportunidades junior \\
\bottomrule
\end{tabular}
\end{center}

\section{Empresas tecnologicas de mineria}

\begin{center}
\begin{tabular}{p{4cm}p{8cm}}
\toprule
\textbf{Empresa} & \textbf{Que hace} \\
\midrule
Komatsu          & FrontRunner (camiones autonomos), AHS, vision \\
Modular Mining   & DISPATCH (despacho con AI), absorbida por Komatsu \\
Caterpillar      & MineStar Command, autonomous trucks \\
Hexagon Mining   & sensores + analytics (ex-Leica) \\
Outotec / Metso  & instrumentacion para flotacion, control \\
Endress+Hauser   & instrumentos calibrados, sensores premium \\
Honeywell Peru   & Experion DCS, integracion \\
Siemens Peru     & TIA Portal, MindSphere IoT \\
ABB Peru         & DCS 800xA, drives \\
Schneider Electric Peru & EcoStruxure, Modicon \\
\bottomrule
\end{tabular}
\end{center}

\section{EPCs e integradores (puerta de entrada para junior)}

\begin{center}
\begin{tabular}{p{4cm}p{8cm}}
\toprule
\textbf{Empresa} & \textbf{Que hace} \\
\midrule
Bechtel        & EPC global, mineria \\
Fluor          & EPC global \\
Hatch          & EPC con foco mining \\
Ausenco        & EPC mid-tier \\
GR Engineering & EPC mid-tier \\
GTD Peru       & integrador chileno con sucursal \\
TGestiona      & integrador (subsidiaria Telefonica) \\
Comesa         & integrador local Lima \\
Intergrupo / Ingelectra & integradores locales \\
\bottomrule
\end{tabular}
\end{center}

% =====================================================================
\chapter{Plataformas de empleo}

\section{Plataformas con ofertas reales}

\begin{center}
\begin{tabular}{p{3cm}p{4cm}p{6cm}}
\toprule
\textbf{Plataforma} & \textbf{URL} & \textbf{Ventaja} \\
\midrule
LinkedIn Jobs       & \texttt{linkedin.com/jobs}    & la mas potente; filtra por empresa, zona, skill \\
Bumeran Peru        & \texttt{bumeran.com.pe}       & portal local lider \\
Computrabajo Peru   & \texttt{computrabajo.com.pe}  & gratis, alto volumen \\
Aptitus             & \texttt{aptitus.com}          & ofertas profesionales \\
Laborum             & \texttt{laborum.pe}           & mid-level / professional \\
\bottomrule
\end{tabular}
\end{center}

\section{Como buscar especificamente}

\subsection{Queries para LinkedIn Jobs}

\begin{itemize}[leftmargin=2em]
  \item \texttt{"data scientist"} + \texttt{"Peru"} + \texttt{"mining"}
  \item \texttt{"computer vision"} + \texttt{"Peru"} + \texttt{"mining"}
  \item \texttt{"automation engineer"} + \texttt{"Peru"}
  \item \texttt{"control systems"} + \texttt{"Antamina"}
  \item \texttt{"data scientist"} + \texttt{"Quellaveco"}
  \item \texttt{"machine learning"} + \texttt{"Cerro Verde"}
  \item \texttt{"automatizacion"} + \texttt{"minera"} + \texttt{"Peru"}
  \item \texttt{"vision artificial"} + \texttt{"industrial"} + \texttt{"Peru"}
\end{itemize}

\subsection{Queries para Bumeran / Computrabajo}

\begin{itemize}[leftmargin=2em]
  \item ingeniero automatizacion mineria
  \item data scientist mina
  \item control e instrumentacion
  \item SCADA programador
  \item ingeniero PLC
  \item vision artificial industrial
  \item analista de datos minero
\end{itemize}

\section{Estrategias avanzadas de busqueda}

\begin{tip}
\textbf{1. Reclutadores directo:} en LinkedIn buscar
\texttt{"talent acquisition Antamina"} o \texttt{"recruiter Cerro Verde"}.
Conectar con un mensaje breve presentandote.

\textbf{2. Programas trainee:} cada mina tiene programa anual (Antamina
"Practicas", Cerro Verde "Programa Talento Joven", Anglo American "Graduate
Programme"). Seguir LinkedIn de las empresas para no perderse las convocatorias.

\textbf{3. Eventos del sector:}
\begin{itemize}[leftmargin=1em]
  \item PERUMIN (Arequipa, bianual) --- la convencion mas grande.
  \item EXTEMIN (asociado a PERUMIN) --- feria de tecnologia.
  \item Mining \& Energy LATAM (Lima, anual).
  \item Capacitaciones IIMP (Instituto de Ingenieros de Minas Peru).
\end{itemize}

\textbf{4. Gerentes de innovacion en LinkedIn:} muchas minas tienen
"Innovation Manager" o "Digital Transformation Lead". Seguirlos. Postean
ofertas de su equipo antes que salgan a portales.
\end{tip}

% =====================================================================
\chapter{Certificaciones que aceleran}

\begin{center}
\begin{tabular}{p{4.5cm}p{4cm}p{2.5cm}p{2cm}}
\toprule
\textbf{Certificacion} & \textbf{Donde se cursa} & \textbf{Costo aprox.} & \textbf{Tiempo} \\
\midrule
Siemens TIA Portal Basic & TecsupVirtual / Siemens Peru & USD 300--500 & 1--2 meses \\
Rockwell RSLogix 5000 Basic & reps Rockwell Peru & USD 400--600 & 1 mes \\
ISA-99 Cybersecurity industrial & online ISA & USD 1000+ & 2--3 meses \\
AVEVA InTouch / Wonderware & AVEVA Peru & USD 500--800 & 1--2 meses \\
AWS Machine Learning   & online (gratis) & USD 0 & 2 meses \\
Coursera Data Engineering & Coursera Plus & USD 50/mes & 4--6 meses \\
Operario PLC nivel 1 (Tecsup) & Tecsup Lima & USD 500 & 1 mes \\
Diplomado Automatizacion (UNI) & UNI / TECSUP & USD 1500--2500 & 6 meses \\
\bottomrule
\end{tabular}
\end{center}

\textbf{Recomendacion priorizada} para alguien que viene del proyecto INTISAT:

\begin{enumerate}[leftmargin=2em]
  \item \textbf{Siemens TIA Portal Basic}: te da el lenguaje del sector.
  \item \textbf{AWS Machine Learning}: capitaliza tu skill en IA.
  \item \textbf{ISA-99 Cybersecurity}: especialidad premium con escasez.
  \item \textbf{Coursera Data Engineering}: te abre roles de data engineer.
\end{enumerate}

% =====================================================================
\chapter{Como vender tu skill set actual}

Tu trabajo en el proyecto CubeSat-EdgeAI tiene \textbf{valor directo en mineria}:

\begin{center}
\begin{tabular}{p{6cm}p{6cm}}
\toprule
\textbf{Tu skill (CubeSat)} & \textbf{Como lo vendes en mineria} \\
\midrule
Pipeline IA con StarDist, OpenCV, PyTorch & "vision por computadora industrial — deteccion de defectos, granulometria, reconocimiento de mineral" \\
Linux embebido + systemd                  & "edge computing en sitios remotos" \\
I2C / SPI / protocolo custom              & "integracion entre sensores y SCADA" \\
Power budget + telemetria                 & "monitoreo IoT industrial" \\
OTA con rollback                           & "deployment seguro en plantas remotas" \\
Autonomia + watchdog                       & "sistemas tolerantes a fallas" \\
Documentacion en LaTeX + GitHub           & "ingenieria sistematica, capaz de redactar reportes profesionales" \\
\bottomrule
\end{tabular}
\end{center}

Con esto te posicionas como:

\begin{itemize}[leftmargin=2em]
  \item \textbf{Junior en mineria} hoy: 4 500--7 000 S/ + bonos.
  \item \textbf{Semi-senior} en 2 años: 7 000--12 000 S/.
  \item \textbf{Senior especialista} en 7 años: 14 000--22 000 S/.
\end{itemize}

% =====================================================================
\chapter{Plan de accion concreto}

\section{Esta semana}

\begin{enumerate}[leftmargin=2em]
  \item Actualizar perfil LinkedIn destacando:
        \begin{itemize}[leftmargin=1em]
          \item Pipeline IA con OpenCV / StarDist / PyTorch
          \item Linux embebido + systemd
          \item Programacion en Python
          \item Proyecto INTISAT con link al GitHub
        \end{itemize}
  \item Configurar alertas en LinkedIn Jobs con las queries del capitulo 9.
  \item Configurar cuenta en Bumeran y Computrabajo con las mismas keywords.
  \item Conectar con 5 personas del rubro (ingenieros de Antamina, Cerro Verde,
        Anglo American, integradores) con mensaje breve presentandote.
\end{enumerate}

\section{Mes 1}

\begin{enumerate}[leftmargin=2em]
  \item Postular a 5--10 ofertas por semana (junior / trainee / practicas).
  \item Inscribirse en \textbf{Siemens TIA Portal Basic} online.
  \item Completar curso \textbf{AWS Machine Learning Foundations} (gratis).
  \item Empezar a leer documentacion de \textbf{Wonderware InTouch} o \textbf{Ignition}.
\end{enumerate}

\section{Mes 2-3}

\begin{enumerate}[leftmargin=2em]
  \item Conseguir \textbf{primera entrevista}.
  \item Si hay convocatoria: participar en \textbf{PERUMIN} o algun evento del sector.
  \item Completar TIA Portal Basic.
  \item Empezar a programar un \textbf{LOGO! Siemens} ($\sim$\$200) o \textbf{CODESYS Runtime} (gratis) en una RPi
        para tener proyecto demostrable.
\end{enumerate}

\section{Mes 4-6}

\begin{enumerate}[leftmargin=2em]
  \item Idealmente: estar contratado en algun rol junior.
  \item Si no: pivotear --- aplicar a EPCs / integradores (mas accesibles).
  \item Continuar con certificacion AVEVA / Wonderware.
  \item Networking sostenido con la comunidad.
\end{enumerate}

\section{Año 2}

\begin{enumerate}[leftmargin=2em]
  \item Consolidar primer empleo, acumular experiencia en sitio si posible.
  \item Considerar ISA-99 Cybersecurity como especialidad premium.
  \item Evaluar postulacion a Australia/Canada como meta de mediano plazo.
\end{enumerate}

% =====================================================================
\chapter{Universidades y formacion en Peru}

\begin{center}
\begin{tabular}{p{6cm}p{7cm}}
\toprule
\textbf{Institucion} & \textbf{Foco} \\
\midrule
UNI (Universidad Nacional de Ingenieria)        & mineria, INTISAT, automatizacion \\
PUCP (Pontificia Universidad Catolica Peru)     & mecatronica, ingenieria informatica \\
UNAS (Universidad Agraria La Molina)             & agroindustria \\
UNSA (San Agustin Arequipa)                      & cerca de Cerro Verde \\
Tecsup (instituto tecnico)                       & cursos cortos especializados \\
TECSUP Diplomados                                & automatizacion, IoT, data \\
\bottomrule
\end{tabular}
\end{center}

\textbf{Tecsup} en particular es muy valorado en industria: programa
"Tecnico en Automatizacion Industrial" + diplomados especializados.
La industria paga estos cursos a sus empleados.

% =====================================================================
\chapter{Comparativa LATAM y migracion}

\begin{center}
\begin{tabular}{lll}
\toprule
\textbf{Pais} & \textbf{Madurez sector} & \textbf{Atractivo} \\
\midrule
Argentina & Alta. Lider regional. Satellogic, Innova Space. & alto si te interesa NewSpace \\
Brasil    & Alta. Industria propia. INPE, Visiona.          & alto, idioma diferente \\
Mexico    & Media. En crecimiento. Dereum Labs.              & medio \\
Chile     & Media. Inversion privada. StarHorse.             & alto, lengua \\
Colombia  & Baja. Solo academico.                            & bajo \\
Peru      & Baja. Solo academico (excepto mineria).          & ref \\
Bolivia / Ecuador & Muy baja.                                  & bajo \\
\bottomrule
\end{tabular}
\end{center}

\textbf{Para mineria especificamente}: Chile y Peru son los polos. Despues
Australia / Canada (donde el sueldo se duplica o triplica).

% =====================================================================
\chapter{Conclusion y reflexion}

\begin{itemize}[leftmargin=2em]
  \item \textbf{Peru NO esta atrasado en industrial}: la mineria moderna usa
        lo mismo que cualquier pais del mundo.
  \item \textbf{Peru SI esta atrasado en aeroespacial}: pero el INTISAT
        es justamente lo que cierra esa brecha.
  \item \textbf{Tu skill set en CubeSat / IA aplica directo en mineria}:
        es un mercado mucho mas grande y mejor pago.
  \item \textbf{Aeroespacial paga peor que mineria} en Peru. Si te interesa
        aeroespacial, hacelo por vocacion. Si necesitas plata, mineria.
  \item \textbf{El proyecto INTISAT te abre dos puertas}: aeroespacial
        (poco mercado en Peru) o vision computacional industrial (mercado
        enorme en mineria).
  \item \textbf{Lo que paga es conectar OT (operational technology) con IT
        (information technology) + AI}. Tu perfil es exactamente eso.
  \item \textbf{El camino tipico desde aqui}: 6 meses para conseguir el
        primer rol junior, 7 años para senior. 15 años para roles
        directivos o consultoria.
\end{itemize}

\begin{tip}
\textbf{Resumen en una linea}: Peru tiene infraestructura industrial
real con sueldos competitivos, especialmente en mineria. Tu proyecto
aeroespacial academico es un \textbf{trampolin tecnico solido} hacia
el sector industrial productivo. El plan de transicion concreto cabe
en 6 meses.
\end{tip}

% =====================================================================
\chapter*{Recursos para profundizar}
\addcontentsline{toc}{chapter}{Recursos para profundizar}

\textbf{Plataformas de empleo (verificar URLs vigentes)}:
\begin{itemize}[leftmargin=2em]
  \item LinkedIn Jobs --- el que mas se usa en sector profesional
  \item Bumeran Peru --- portal local
  \item Computrabajo Peru --- gratis, alto volumen
  \item Aptitus --- profesional / mando medio
\end{itemize}

\textbf{Asociaciones del sector}:
\begin{itemize}[leftmargin=2em]
  \item IIMP (Instituto de Ingenieros de Minas del Peru)
  \item ISA Peru Section (International Society of Automation)
  \item ICONTEC Peru (normas tecnicas)
\end{itemize}

\textbf{Eventos anuales / bianuales}:
\begin{itemize}[leftmargin=2em]
  \item PERUMIN (Arequipa, bianual)
  \item EXTEMIN (con PERUMIN)
  \item Mining \& Energy LATAM (Lima)
  \item Capacitaciones IIMP (varias por año)
  \item Expo MIN (Lima)
\end{itemize}

\textbf{Lectura recomendada}:
\begin{itemize}[leftmargin=2em]
  \item Norma ISA-95 (modelo piramide automatizacion)
  \item Norma ISA-99 / IEC 62443 (cyberseguridad industrial)
  \item Norma IEC 61131-3 (lenguajes PLC)
  \item Norma MIL-STD-810 (tests ambientales)
\end{itemize}

\end{document}
